\chapter{Consumers and Preferences}
\label{chap:preferences}

\section{Learning objectives}

After completing this chapter, you should be able to:
\begin{itemize}
  \item describe a consumer's choice problem using consumption bundles;
  \item use the symbols $\succeq$, $\succ$, and $\sim$ correctly;
  \item state and interpret the principal axioms of rational preferences;
  \item explain monotonicity, convexity, and their graphical implications; and
  \item recognise both the usefulness and the limits of the rational-choice model.
\end{itemize}

\section{Choice theory}

Consumer theory begins with a simple question: how does a person choose among
the alternatives available to them? The answer will eventually depend on
income and prices, but first we need a way to describe what the person likes
or values. Choosing between two lunch options, two mobile plans, or two ways
to spend a free evening all begins with this ranking. This is the role of
\keyterm{preferences}.

The standard model does not assume that every consumer likes the same things,
nor that consumers care only about material goods. A person may value leisure,
friendship, health, environmental quality, status, or ethical consumption.
What the model assumes is more modest: when faced with alternatives, the
consumer can compare them consistently enough to make a choice.

\section{Consumption bundles}

Suppose that a consumer chooses quantities of two goods, $x$ and $y$. A
\keyterm{consumption bundle} is an ordered pair $(x,y)$ describing the
quantities consumed. For example, $(3,5)$ may mean three units of food and
five units of clothing. The order matters: $(3,5)$ is not the same bundle as
$(5,3)$ when the two goods are different.

Using two goods is a simplification, but it is powerful. It lets us draw
choices on a diagram and focus on the trade-off between alternatives. Later,
$x$ and $y$ can stand for broad categories such as present consumption and
future consumption, private goods and environmental quality, or leisure and
income.

\section{Describing preferences}

Let $A$ and $B$ denote two bundles. The notation
\[
A\succeq B
\]
means that the consumer considers $A$ \term{at least as good as} $B$. It includes two possibilities: the consumer may strictly prefer $A$, or may be indifferent between $A$ and $B$.

The notation
\[
A\succ B
\]
means that $A$ is \keyterm{strictly preferred} to $B$. The consumer chooses $A$ over $B$ when both are available. Finally,
\[
A\sim B
\]
means that the consumer is \keyterm{indifferent} between the two bundles: each is considered at least as good as the other.

The strict and indifferent relations can be defined from the weak-preference
relation:
\[
A\succ B
\quad\Longleftrightarrow\quad
A\succeq B \text{ and } B\not\succeq A,
\]
and
\[
A\sim B
\quad\Longleftrightarrow\quad
A\succeq B \text{ and } B\succeq A.
\]
With complete preferences, exactly one of $A\succ B$, $B\succ A$, or
$A\sim B$ holds.

Preferences rank alternatives. They do not yet explain \emph{why} one bundle
is preferred. The reasons may be personal, cultural, moral, or social. For
example, two people can rank the same vegetarian meal differently for reasons
of taste, health, cost, or ethics. The economic model works with the ranking
itself.

\begin{keyidea}
\textbf{Three ideas that must not be confused.}
\begin{itemize}
  \item \textbf{Preferences} rank bundles: \emph{I consider $A$ at least as good as $B$.}
  \item A \textbf{feasible set} lists the bundles that are actually available.
  \item \textbf{Choice} selects a most-preferred feasible bundle.
\end{itemize}
A person may prefer one bundle but choose another because the preferred bundle
is unaffordable, unavailable, or ruled out by time or institutional
constraints.
\end{keyidea}

\begin{exercisebox}
\textbf{Running example: the campus-caf\'e menu.} Laura ranks weekly bundles
of coffees and lunches. Strong monotonicity tells us that $(4,5)$ is preferred
to $(3,4)$, since it contains more of both goods. It does not tell us whether
she prefers $(5,3)$ or $(3,5)$: that depends on her own preferences. If the
caf\'e runs out of lunches, the feasible set changes even if Laura's
preferences do not.
\end{exercisebox}

\section{Rational preferences}

In introductory microeconomics, preferences are called \keyterm{rational} if
they satisfy a set of consistency properties. In this technical use, rational
means internally coherent. It does not mean selfish, perfectly informed,
emotionally detached, morally admirable, or always correct. The label is
about the pattern of choices, not about whether an observer approves of them.

\subsection{Completeness}

Preferences are \keyterm{complete} if, for any two bundles $A$ and $B$, the
consumer can compare them: either $A\succeq B$ or $B\succeq A$, or both. The
consumer need not have a strong opinion, but complete preferences rule out
being unable to compare the options at all.

Completeness is a useful approximation. In real life, people may face unfamiliar alternatives or may postpone a decision. The assumption becomes more plausible when the alternatives are well understood and the consumer has enough time to decide.

\subsection{Transitivity}

Preferences are \keyterm{transitive} if
\[
A\succeq B \quad \text{and} \quad B\succeq C
\quad \Longrightarrow \quad A\succeq C.
\]
Transitivity rules out cycles of preference. If a consumer preferred $A$ to
$B$, $B$ to $C$, and $C$ to $A$, a clever seller could repeatedly exchange
bundles and leave the consumer worse off. Transitivity therefore protects the
idea that choice is consistent over comparable alternatives. It is the formal
version of not changing your mind in a circle.

\subsection{Reflexivity}

Preferences are \keyterm{reflexive} if every bundle is at least as good as itself:
\[
A\succeq A.
\]
Under the usual definition of completeness, reflexivity follows immediately by
comparing a bundle with itself. It is stated explicitly because it is useful
to keep the weak-preference relation formally clear.

\section{More is better: monotonicity}

The \keyterm{strong monotonicity} assumption, often described informally as
\emph{more is better}, states that more of a desirable good is preferred to
less, holding everything else constant. If $A=(x_A,y_A)$ and $B=(x_B,y_B)$,
with $x_A\geq x_B$ and $y_A\geq y_B$, and at least one inequality strict,
then strong monotonicity implies $A\succ B$.

\begin{figure}[!ht]
\centering
\begin{tikzpicture}[x=.62cm,y=.62cm]
  \fill[PPEGold!15] (0,0) rectangle (4,4);
  \fill[PPEAccentLight] (4,4) rectangle (7.1,7.1);
  \draw[->,PPEInk] (0,0) -- (7.45,0) node[right] {good $x$};
  \draw[->,PPEInk] (0,0) -- (0,7.45) node[above] {good $y$};
  \draw[PPEMuted,dashed] (4,0) -- (4,7.1);
  \draw[PPEMuted,dashed] (0,4) -- (7.1,4);
  \fill[PPENavarraRed] (4,4) circle (1.7pt);
  \node[below left,PPENavarraRed] at (4,4) {$A$};
  \node[align=center,font=\sffamily\scriptsize] at (5.65,5.65)
    {strictly preferred\\to $A$};
  \node[align=center,font=\sffamily\scriptsize] at (2,2)
    {strictly worse\\than $A$};
  \node[align=center,font=\sffamily\scriptsize] at (2,5.65)
    {not ranked from\\monotonicity alone};
  \node[align=center,font=\sffamily\scriptsize] at (5.65,2)
    {not ranked from\\monotonicity alone};
\end{tikzpicture}
\caption{What strong monotonicity can rank relative to a reference bundle
$A$. The north-west and south-east regions require additional information
about preferences.}
\label{fig:monotonicity-regions}
\end{figure}

For ordinary goods, this is often a reasonable starting point. It implies
that a consumer will not voluntarily throw away a desirable good. However,
some goods can become ``bads'' beyond a certain point: pollution, noise, risk,
or even food after satiation. A second dessert may be welcome; a tenth may
not be. The assumption must therefore be matched to the question. We can
define a good so that larger quantities are desirable within the relevant
range, or use a model that explicitly allows satiation.

\section{Continuity and convexity}

\keyterm{Continuity} rules out arbitrary jumps in preference rankings.
Informally, if $A\succ B$, sufficiently small changes to either bundle should
not reverse that strict ranking. If you clearly prefer one weekly food plan to
another, changing it by one piece of fruit should not normally turn it into
your least preferred option. Continuity is a regularity assumption that,
together with completeness, transitivity, and standard assumptions about the
consumption set, helps ensure that preferences can be represented by a
continuous utility function.

\keyterm{Convexity} often captures a preference for variety or balanced
bundles. It says that a mix can be at least as attractive as an extreme. Let
\[
\lambda A+(1-\lambda)B
=
\bigl(\lambda x_A+(1-\lambda)x_B,\,
      \lambda y_A+(1-\lambda)y_B\bigr),
\qquad 0\leq\lambda\leq1.
\]
Preferences are convex if, whenever $A$ and $B$ are each at least as good as
$C$, every mixture of $A$ and $B$ is also at least as good as $C$. In
particular, if a consumer is indifferent between $A$ and $B$, a mixture is at
least as good as either endpoint. With \keyterm{strict convexity}, a genuine
mixture of two distinct indifferent bundles is strictly preferred to either
extreme.

\begin{figure}[!ht]
\centering
\begin{tikzpicture}[x=.58cm,y=.58cm]
  \draw[->,PPEInk] (0,0) -- (8.25,0) node[right] {good $x$};
  \draw[->,PPEInk] (0,0) -- (0,8.25) node[above] {good $y$};
  \draw[PPEGold,line width=.95pt] (1,7) -- (7,1);
  \fill[PPENavarraRed] (1,7) circle (1.7pt);
  \fill[PPENavarraRed] (7,1) circle (1.7pt);
  \fill[PPEBurgundy] (4,4) circle (1.9pt);
  \node[above right,PPENavarraRed] at (1,7) {$A=(1,7)$};
  \node[below left,PPENavarraRed] at (7,1) {$B=(7,1)$};
  \node[above right,PPEBurgundy,align=left,font=\sffamily\scriptsize] at (4,4)
    {$M=(4,4)$\\$M\succeq A\sim B$};
\end{tikzpicture}
\caption{Convexity and mixtures. If $A\sim B$, their midpoint $M$ is at
least as good as either endpoint; strict convexity makes the preference
strict.}
\label{fig:convexity-mixture}
\end{figure}

For example, a student may prefer a balanced allocation of meals and leisure
to an extreme allocation that gives only one of them. This interpretation
assumes that the goods are divisible, or that a mixture can be understood as
an average across time. Convexity allows perfect substitutes, whose
indifference curves are straight lines. Strict convexity gives the familiar
strictly bowed-in indifference curves introduced in the next chapter.

\section{From preferences to choice}

Preferences alone do not determine an observed choice. A consumer may prefer
a luxury holiday to a local holiday but lack the income to buy it. Choice
requires both preferences and a \keyterm{feasible set}: the alternatives that
are actually available. This distinction helps us avoid reading too much into
what people buy. In the next chapter, utility functions and indifference
curves will represent preferences; after that, budget constraints will
describe the feasible set.

The model is deliberately modest. It describes what a person chooses when
their preferences and constraints are given. It does not claim that
preferences are fixed forever or unaffected by social context. Advertising,
habit, family, institutions, social norms, and political power can shape the
alternatives people perceive and the preferences they express.

\begin{keyidea}
\textbf{The sequence of consumer theory.}
\[
\text{preferences}
\longrightarrow
\text{utility representation}
\longrightarrow
\text{indifference curves}
\longrightarrow
\text{constrained choice}.
\]
This chapter establishes the first step only.
\end{keyidea}

\section*{Chapter Review}

\subsection*{Key concepts}
\begin{itemize}
  \item \textbf{Consumption bundle}: a list of quantities of goods consumed.
  \item \textbf{Preference relation}: a ranking of alternative bundles.
  \item \textbf{Completeness}: every pair of bundles can be compared.
  \item \textbf{Transitivity}: preference rankings do not form cycles.
  \item \textbf{Strong monotonicity}: more of a desirable good is preferred to less, other things equal.
  \item \textbf{Continuity}: small changes do not cause arbitrary jumps in preferences.
  \item \textbf{Convexity}: mixtures of equally valued bundles are at least as good as extremes.
  \item \textbf{Feasible set}: the alternatives available to the consumer.
\end{itemize}

\subsection*{Notation}
\begin{center}
\begin{tabular}{ll}
\toprule
Notation & Meaning \\
\midrule
$A\succeq B$ & $A$ is at least as good as $B$ \\
$A\succ B$ & $A$ is strictly preferred to $B$ \\
$A\sim B$ & The consumer is indifferent between $A$ and $B$ \\
\bottomrule
\end{tabular}
\end{center}

\section*{Exercises}

\subsection*{Conceptual questions}
\begin{enumerate}[label=\arabic*.]
  \item Define a consumption bundle. Why are two-good models useful even though consumers purchase many goods?
  \item Explain the difference between $A\succeq B$, $A\succ B$, and $A\sim B$.
  \item State completeness and transitivity precisely. Give one example of a situation in which each assumption may be a useful approximation.
  \item Does rationality in consumer theory mean selfishness? Explain.
  \item Explain why monotonicity may be inappropriate for pollution or risk.
\end{enumerate}

\subsection*{Graphical exercises}
\begin{enumerate}[resume,label=\arabic*.]
  \item In a diagram with good $x$ on the horizontal axis and good $y$ on the vertical axis, place the bundles $A=(2,2)$, $B=(4,2)$, $C=(3,5)$, and $D=(6,6)$. Assuming strong monotonicity, compare $A$ with $B$, $A$ with $C$, and $A$ with $D$. Can you compare $B$ and $C$ without more information?
  \item Let $A=(4,0)$ and $B=(0,4)$, and suppose $A\sim B$. Plot both bundles and their midpoint $M=(2,2)$. What does convexity imply about $M$? What additional conclusion follows under strict convexity?
  \item Draw a case in which strong monotonicity holds and identify the direction in which preferred bundles lie.
\end{enumerate}

\subsection*{Analytical exercises}
\begin{enumerate}[resume,label=\arabic*.]
  \item Suppose $A\succ B$, $B\sim C$, and $C\succeq D$. Show that $A\succ D$.
  \item Holding the decision context fixed, a consumer has the following rankings: $A\succ B$, $B\succ C$, and $C\succ A$. Does this consumer satisfy transitivity? Explain the practical problem caused by this pattern.
  \item Let $A=(4,3)$ and $B=(2,3)$. Under strong monotonicity, write the relevant preference relation and justify it.
  \item Suppose the feasible set is $F=\{B,C\}$. If $A\succ B\succ C$, but $A\notin F$, which bundle is chosen? Explain why preferences alone do not determine observed choice.
  \item Suppose $A\sim B$ and $B\sim C$. Use the definition of indifference and transitivity to show that $A\sim C$.
\end{enumerate}

\subsection*{Exam-style multiple-choice questions}
\begin{enumerate}[resume,label=\arabic*.]
  \item If $A\sim B$, then: (a) $A\succ B$; (b) $B\succ A$; (c) $A\succeq B$ and $B\succeq A$; or (d) the consumer cannot compare $A$ and $B$.
  \item Transitivity means that if $A\succeq B$ and $B\succeq C$, then: (a) $C\succeq A$; (b) $A\succeq C$; (c) $A\sim C$; or (d) $B\succ A$.
  \item Convex preferences imply that, when a consumer is indifferent between two bundles: (a) only extreme bundles can be chosen; (b) a mixture is at least as good as either bundle; (c) more is always worse; or (d) the consumer has no preferences.
\end{enumerate}

\subsection*{PPE discussion question}

\noindent Consumer theory treats preferences as given. Discuss whether this is an acceptable simplification for public policy. In your answer, consider advertising, education, social norms, and the possibility that preferences themselves may be shaped by institutions.
